\section{BSGS 及其拓展}

\subsection{BSGS 算法}

\verb|BSGS(a,b,p)| 求 $a^x\equiv b\pmod{p}$ 的最小非负整数解，无解返回 $-1$，复杂度 $\mathcal{O}(\sqrt{p})$。要求 $\gcd(a,p)=1$（$p$ 为素数时自然满足）；$a,p$ 不互质或 $p$ 为合数时改用 \verb|exBSGS|。

\begin{lstlisting}
using i64 = long long;
i64 BSGS(i64 a,i64 b,const i64& p){ // a ^ x = b \pmod p
    a %= p,b %= p;
    if (b == 1) return 0;
    const i64 t = sqrt(p);
    std::unordered_map<long long,long long> mp;
    i64 ap = 1;
    for (i64 i = 0,res = b;i < t;i++){
        mp[res] = i;
        res = (i64)((__int128)res * a % p);
        ap = (i64)((__int128)ap * a % p);
    }
    for (i64 i = t,v = ap;i < p + t;i += t,v = (i64)((__int128)v * ap % p)){
        if (mp.contains(v)) {
            return i - mp[v];
        }
    }
    return -1;
}
\end{lstlisting}

\subsection{exBSGS 算法}

\verb|exBSGS(a,b,p)| 求 $a^x\equiv b\pmod{p}$ 的最小非负整数解，无解返回 $-1$，复杂度约为 $\mathcal{O}(\sqrt p)$。适用于 $a,p$ 不互质、$p$ 为合数等 \verb|BSGS| 失效的情形。

前置模板：\textbf{\texttt{拓展欧几里得算法}} 与上式 \textbf{\texttt{BSGS}}（递归基直接调用）。

\begin{lstlisting}
using i64 = long long;
i64 exBSGS(i64 a,i64 b,const i64& p){
    a %= p,b %= p;
    if (b % p == 1 % p) return 0;
    i64 d = std::gcd(a,p);
    if (b % d != 0) return -1;
    if (d > 1) {
        i64 x,y;
        i64 np = p / d;
        exgcd(a / d,np,x,y);
        i64 nb = (i64)(((__int128)(b / d) * x % np + np) % np);
        auto ans = exBSGS(a,nb,np);
        return ans == -1 ? -1 : ans + 1;
    }
    return BSGS(a,b,p);
}
\end{lstlisting}